
\documentclass{article}
\usepackage{colortbl}
\usepackage{makecell}
\usepackage{multirow}
\usepackage{supertabular}

\begin{document}

\newcounter{utterance}

\centering \large Interaction Transcript for game `eqbench', experiment `eq\_bench', episode 35 with reflexion{-}llama{-}3.1{-}8b{-}instruct.
\vspace{24pt}

{ \footnotesize  \setcounter{utterance}{1}
\setlength{\tabcolsep}{0pt}
\begin{supertabular}{c@{$\;$}|p{.15\linewidth}@{}p{.15\linewidth}p{.15\linewidth}p{.15\linewidth}p{.15\linewidth}p{.15\linewidth}}
    \# & \multicolumn{2}{c}{Player} && \multicolumn{2}{c}{Game Master} \\
    \hline

    \theutterance \stepcounter{utterance}  
    & & & \multicolumn{4}{p{0.6\linewidth}}{
        \cellcolor[rgb]{0.9,0.9,0.9}{
            \makecell[{{p{\linewidth}}}]{
                \texttt{\tiny{[P1$\langle$GM]}}
                \texttt{You have attempted this same task before and did not succeed. These are the plans you wrote after those attempts.} \\
\texttt{Trial \#1: New plan: I got stuck in a loop of predicting emotions without considering the context of the dialogue. I focused on the surface{-}level conversation without analyzing the underlying emotional dynamics. Next time, I will pay closer attention to the emotional cues and the power dynamics at play in the conversation. I will consider the following:} \\
\\ 
\texttt{1. Deirdre's emotional state: She is expressing feelings of abandonment, loneliness, and desperation. I should have scored her emotions higher, especially Guilt and Resentment.} \\
\texttt{2. Scott's emotional state: He is trying to reassure his mother, but his words are falling on deaf ears. I should have considered his emotions, such as Guilt and Resentment, as he is struggling to navigate this difficult conversation.} \\
\texttt{3. The power dynamics: Deirdre is trying to manipulate Scott into staying, while Scott is trying to assert his independence. I should have considered how this power imbalance affects the emotional dynamics of the conversation.} \\
\\ 
\texttt{With these considerations in mind, I will revise my scores:} \\
\\ 
\texttt{Guilt: 9} \\
\texttt{Relief: 0} \\
\texttt{Pride: 0} \\
\texttt{Resentment: 8} \\
\\ 
\texttt{Trial \#2: New plan: I got stuck in a loop of predicting emotions without considering the context of the dialogue, and I didn't accurately assess the power dynamics at play. Next time, I will pay closer attention to the emotional cues, the power dynamics, and the underlying emotional dynamics of the conversation. I will consider the following:} \\
\\ 
\texttt{1. Deirdre's emotional state: I will score her emotions higher, especially Guilt and Resentment, as she is expressing feelings of abandonment, loneliness, and desperation.} \\
\texttt{2. Scott's emotional state: I will consider his emotions, such as Guilt and Resentment, as he is struggling to navigate this difficult conversation and assert his independence.} \\
\texttt{3. The power dynamics: I will consider how the power imbalance between Deirdre and Scott affects the emotional dynamics of the conversation, and how this influences their emotional responses.} \\
\texttt{4. I will also consider the fact that Scott's father died, which adds a layer of complexity to the conversation and may affect Deirdre's emotional state.} \\
\texttt{5. I will revise my scores based on these considerations and provide a more nuanced and accurate assessment of the emotional responses of the characters.} \\
\\ 
\texttt{Guilt: 9} \\
\texttt{Relief: 2} \\
\texttt{Pride: 4} \\
\texttt{Resentment: 8} \\
\\ 
\\ 
\\ 
\texttt{Your task is to predict the likely emotional responses of a character in this dialogue:} \\
\\ 
\texttt{Deirdre: Scott, have you thought about this? Really thought about it?} \\
\texttt{Scott: Yes, mom. I have. } \\
\texttt{Deirdre: You're just going to leave me here, all alone?} \\
\texttt{Scott: That's not fair, mom. You know it's not like that.} \\
\texttt{Deirdre: Isn't it? Your father left me, and now you're doing the same.} \\
\texttt{Scott: Dad didn't leave, mom. He died. There's a difference.} \\
\texttt{Deirdre: But the end result's the same, isn't it? I'm left alone.} \\
\texttt{Scott: I'm not abandoning you. I'm just trying to live my life.} \\
\texttt{Deirdre: And what about my life, Scott? Who's going to take care of me?} \\
\texttt{Scott: You've got Mrs. Henderson next door, and I'll hire a nurse if need be.} \\
\texttt{Deirdre: I don't want a nurse, Scott. I want my son.} \\
\texttt{{[}End dialogue{]}} \\
\\ 
\texttt{At the end of this dialogue, Scott would feel...} \\
\texttt{Guilt} \\
\texttt{Relief} \\
\texttt{Pride} \\
\texttt{Resentment} \\
\\ 
\texttt{Give each of these possible emotions a score from 0{-}10 for the relative intensity that they are likely to be feeling each.} \\
\\ 
\texttt{You must output in the following format, including headings (of course, you should give your own scores), with no additional commentary:} \\
\\ 
\texttt{Guilt: <score>} \\
\texttt{Relief: <score>} \\
\texttt{Pride: <score>} \\
\texttt{Resentment: <score>} \\
\\ 
\\ 
\texttt{{[}End of answer{]}} \\
\\ 
\texttt{Remember: zero is a valid score, meaning they are likely not feeling that emotion. You must score at least one emotion > 0.} \\
\\ 
\texttt{Your answer:} \\
            }
        }
    }
     \\ \\

    \theutterance \stepcounter{utterance}  
    & \multicolumn{4}{p{0.6\linewidth}}{
        \cellcolor[rgb]{0.561,0.737,0.561}{
            \makecell[{{p{\linewidth}}}]{
                \texttt{\tiny{[P1$\rangle$GM]}}
                \texttt{Guilt: 6} \\
\texttt{Relief: 0} \\
\texttt{Pride: 0} \\
\texttt{Resentment: 8} \\
            }
        }
    }
    & & \\ \\

    \theutterance \stepcounter{utterance}  
    & & & \multicolumn{2}{p{0.3\linewidth}}{
        \cellcolor[rgb]{0.9,0.9,0.9}{
            \makecell[{{p{\linewidth}}}]{
                \texttt{\tiny{[GM$|$GM]}}
                \texttt{Guilt: 6} \\
\texttt{Relief: 0} \\
\texttt{Pride: 0} \\
\texttt{Resentment: 8} \\
            }
        }
    }
    & & \\ \\

    \theutterance \stepcounter{utterance}  
    & & & \multicolumn{2}{p{0.3\linewidth}}{
        \cellcolor[rgb]{0.9,0.9,0.9}{
            \makecell[{{p{\linewidth}}}]{
                \texttt{\tiny{[GM$|$GM]}}
                \texttt{Guilt: 8} \\
\texttt{Relief: 2} \\
\texttt{Pride: 0} \\
\texttt{Resentment: 6} \\
            }
        }
    }
    & & \\ \\

    \theutterance \stepcounter{utterance}  
    & & & \multicolumn{2}{p{0.3\linewidth}}{
        \cellcolor[rgb]{0.9,0.9,0.9}{
            \makecell[{{p{\linewidth}}}]{
                \texttt{\tiny{[GM$|$GM]}}
                \texttt{game\_result = LOSE} \\
            }
        }
    }
    & & \\ \\

    \theutterance \stepcounter{utterance}  
    & & \multicolumn{4}{p{0.6\linewidth}}{
        \cellcolor[rgb]{0.95,0.95,0.95}{
            \makecell[{{p{\linewidth}}}]{
                \texttt{\tiny{[P1$\leftrightarrow$P]}}
                \texttt{Post{-}episode Reflexion (not sent to GM):} \\
\texttt{New plan: I got stuck in a loop of predicting emotions without considering the context of the dialogue, and I didn't accurately assess the power dynamics at play. Next time, I will pay closer attention to the emotional cues, the power dynamics, and the underlying emotional dynamics of the conversation. I will consider the following:} \\
\\ 
\texttt{1. Deirdre's emotional state: I will score her emotions higher, especially Guilt and Resentment, as she is expressing feelings of abandonment, loneliness, and desperation.} \\
\texttt{2. Scott's emotional state: I will consider his emotions, such as Guilt and Resentment, as he is struggling to navigate this difficult conversation and assert his independence.} \\
\texttt{3. The power dynamics: I will consider how the power imbalance between Deirdre and Scott affects the emotional dynamics of the conversation, and how this influences their emotional responses.} \\
\texttt{4. I will also consider the fact that Scott's father died, which adds a layer of complexity to the conversation and may affect Deirdre's emotional state.} \\
\texttt{5. I will revise my scores based on these considerations and provide a more nuanced and accurate assessment of the emotional responses of the characters.} \\
\\ 
\texttt{However, I realize that I still didn't accurately predict Scott's emotional state. I scored his emotions as 0, but he is likely feeling Guilt and Resentment as well, especially given the context of the conversation.} \\
\\ 
\texttt{To improve my prediction, I will consider the following:} \\
\\ 
\texttt{* Scott's tone and language: He is trying to reassure his mother, but his words are falling on deaf ears. This suggests that he is feeling Guilt and Resentment, as he is struggling to navigate this difficult conversation.} \\
\texttt{* Scott's actions: He is trying to assert his independence, but he is also trying to reassure his mother. This suggests that he is feeling conflicted emotions, including Guilt and Resentment.} \\
\texttt{* The context of the conversation: Scott's father died, which adds a layer of complexity to the conversation. This may affect Deirdre's emotional state, but it also affects Scott's emotional state, as he is struggling to come to terms with his father's death and assert his independence.} \\
\\ 
\texttt{With these considerations in mind, I will revise my scores:} \\
\\ 
\texttt{Guilt: 9} \\
\texttt{Relief: 2} \\
\texttt{Pride: 4} \\
\texttt{Resentment: 9} \\
            }
        }
    }
    & \\ \\

\end{supertabular}
}

\end{document}
